\documentclass[11pt,a4paper]{article}
\usepackage[utf8]{inputenc}
\usepackage[T1]{fontenc}
\usepackage{amsmath,amssymb,amsthm}
\usepackage[margin=2.4cm]{geometry}
\usepackage{booktabs}
\usepackage{array}
\usepackage{hyperref}

\hypersetup{colorlinks=true,linkcolor=blue,citecolor=blue,urlcolor=blue}

\newcommand{\dd}{\mathrm{d}}
\newcommand{\ii}{\mathrm{i}}
\newcommand{\ee}{\mathrm{e}}
\newcommand{\hb}{\hbar}
\newcommand{\sgn}{\operatorname{sgn}}
\newcommand{\Rea}{\operatorname{Re}}
\newcommand{\Ima}{\operatorname{Im}}
\newcommand{\Li}{\operatorname{Li}}

\title{\textbf{Phase 2g: One-Loop Soliton--Defect Transmission Phase,\\
Closed-Form Defect Stability Spectrum, and KP-Triangle Fusion\\
for the Good Boussinesq Equation}}
\author{Project write-up (Route E, step 2) --- PI: Gustav Delius}
\date{April 2026}

\begin{document}
\maketitle

\begin{abstract}
We compute the one-loop ($O(\hb^0)$) soliton--defect transmission phase for the
good Boussinesq equation $u_{tt} = u_{xx} - u_{xxxx} - 3(u^2)_{xx}$ on the
Corrigan--Zambon (CZ) integrable type-I defect (Route~E, step~2 of
\texttt{report.pdf}). Every nontrivial claim is labelled \textbf{(a)} proved in
the literature, \textbf{(b)} conjectured/heuristic in the literature,
\textbf{(c)} derived and verified in this project, or \textbf{(d)} a guess.

Our main results \textbf{(c)} are:
\begin{enumerate}
\item \textbf{Exact time-dependent phonon scattering on $N$-soliton + defect backgrounds.}
Extending the CZ $1$- and $2$-soliton defect solutions to $3$ solitons and
linearising in an oscillatory wave $E_3 = \varepsilon\,\ee^{\ii(qx-\omega(q)t)}$
($\omega(q) = |q|\sqrt{1+q^2}$) produces an exact time-dependent phonon
scattering eigenmode satisfying all four linearised CZ sewing conditions at
$x=0$ for all $t \in \mathbb{R}$, with zero reflection ($R(q)=0$), zero
asymptotic evanescent tails, and factorised transmission amplitude
$T_{\mathrm{tot}}(q) = \big(\prod_{j=1}^N A(j,\mathrm{ph}(q))\big)z_{\mathrm{ph}}(q;q_D,\eta)$.
\item \textbf{KP-vertex factorisation and 2-root dipole decomposition.}
In terms of the single-KP-root phonon factor
$\mathcal{Z}(q;p) := (p_{\mathrm{ph},b}(q)-p)/(p_{\mathrm{ph},a}(q)-p)$, a
defect at KP vertex $q_D = -s/(2\sqrt{3})$ ($s=\eta\sigma$) has
$z_{\mathrm{ph}}(q;q_D,\eta) = \mathcal{Z}(q;q_D)^\eta$, while a soliton on edge
$(p_a,p_b)$ has $A(\mathrm{sol}(p_a,p_b),\mathrm{ph}(q)) = \mathcal{Z}(q;p_b)/\mathcal{Z}(q;p_a)$.
Setting $q_D = -(\cos\gamma)/\sqrt{3}$ ($\gamma = \arccos(s/2)$), the defect
phase-shift derivative decomposes into a \emph{2-root $(-\eta,+\eta)$ dipole}
$\theta_D'(q;s,\eta) = \Rea\sum_{s'=1}^2 c_{D,s'}\Theta'(q;\psi_{s'})$ over the
universal elementary kernel $\Theta'(q;\psi)$ of \texttt{oneloop\_phase.pdf},
where $\sin\psi_1 = q_+ - q_D$ and $\sin\psi_2 = q_D - q_-$ are the two root
differences between $q_D$ and the remaining two vertices $q_\pm$ of its KP
triple. Subtracting the two defect dipoles at $p_b$ and $p_a$ yields
$(+1,-1,0) - (0,+1,-1) = (+1,-2,+1)$, reproducing the 3-root soliton kernel
$\theta_S'(q;\alpha,\epsilon)$.
\item \textbf{Closed-form algebraic solution of the defect vacuum growing-mode spectrum.}
The $4\times 4$ linearised sewing determinant for localised modes
$\omega = \ii\gamma_{\mathrm{grow}}$ factors in $S = \kappa_1+\kappa_2 = \sqrt{1+2\gamma_{\mathrm{grow}}}$
into two quadratics $P_\mp(S;\eta\sigma)$ with perfect-square discriminants
$\big(\frac{\sqrt{3}\pm 1}{2}\big)^2(4-\sigma^2)$, giving the exact eigenvalues
\[
  \omega_{1,2}(\sigma,\eta) = \ii\,\frac{\eta\sigma\sqrt{4-\sigma^2} \pm \sqrt{3}(2-\sigma^2)}{8},
\]
which explain analytically the thresholds $s=-1$ (where $\omega_1$ emerges from
$\omega=0$ with a $-2\pi$ Levinson jump), $s=+1$ (where $\omega_2$ emerges from
$\omega=0$ with a second $-2\pi$ Levinson jump), and $s=+2$ (where $\omega_1,\omega_2$
collide at $\ii\sqrt{3}/4$ and split into a complex-conjugate oscillatory pair).
\item \textbf{Closed-form $3\times 2$ dilogarithm formula, $(E,P)$ absorption balance, and KP-triangle closure.}
Because the static defect has no translational zero mode ($\theta_D^{\mathrm{phys}}(0^\pm)=0$),
the raw and antisymmetrised Faddeev--Korepin integrals coincide ($I_{SD} = J_{SD}^{\mathrm{phys}}$)
and evaluate in closed form as a $3\times 2$ sum of Legendre $\chi_2$ functions.
Across every soliton absorption process $D(q_i)+S_i \to D(q_{i+1})$ ($S_i=\{q_i,q_{i+1}\}$),
the one-loop defect vacuum energy and momentum shifts satisfy exact balance
$(\Delta E_1^{\mathrm{ren}}(S_i),\Delta P_1(S_i)) = \rho_i\big(\Delta E_D^{\mathrm{ren}}(q_{i+1})-\Delta E_D^{\mathrm{ren}}(q_i),\,\Delta P_D(q_{i+1})-\Delta P_D(q_i)\big)$,
and the one-loop phases obey exact defect fusion
$I_{SD}(3;q_{i+1},\eta)-I_{SD}(3;q_i,\eta) = \rho_i\,I_{SS}(3,S_i)$ with the
intrinsic absorption sign $\rho_i = \eta\sgn(q_i-q_{i+1}) \in \{+1,-1\}$
($\rho_{\mathrm{tree}} = \epsilon_3\rho_i$). Consequently, around every closed KP
triangle $q_1\to q_2\to q_3\to q_1$, the raw Faddeev--Korepin one-loop two-soliton
phase $I_{SS}$ is an exact coboundary ($\sum_{i=1}^3\rho_i\,I_{SS}(3,S_i)=0$),
whereas adding the $+\frac{\pi}{2}\delta_{\epsilon_3,\epsilon_{S_i}}$ zero-mode
boundary constant of $J_{SS}$ (or the $\pm\frac{\pi}{2}$ jump of the $\Gamma$-ansatz)
leaves a non-zero coboundary defect $\frac{\pi}{2}\sum_{i=1}^3\rho_i\,\delta_{\epsilon_3,\epsilon_{S_i}} \in \{0,\pm\frac{\pi}{2}\}$.
\end{enumerate}
\end{abstract}

\tableofcontents

% ==============================================================================
\section{Setup, conventions, and the Route E question at one loop}
\label{sec:setup}
% ==============================================================================

We work in the fixed project conventions (\texttt{README.md},
\texttt{defect\_semiclassics.pdf}, \texttt{oneloop\_phase.pdf}) for the
\textbf{good} Boussinesq equation \textbf{(a)}:
\begin{equation}\label{eq:good_bq}
  u_{tt} = u_{xx} - u_{xxxx} - 3(u^2)_{xx}, \qquad
  u = \varphi_x = 2(\ln\tau)_{xx},
\end{equation}
with bulk Lagrangian density and canonical symplectic form
\begin{equation}\label{eq:lagrangian}
  \mathcal{L} = \frac{1}{2}\varphi_t^2 - \frac{1}{2}\varphi_x^2
  + \varphi_x^3 - \frac{1}{2}\varphi_{xx}^2, \qquad
  \Omega = \int_{-\infty}^\infty \dd\pi\wedge\dd\varphi, \quad \pi = \varphi_t,
\end{equation}
in which $\hb$ is normalised. The Corrigan--Zambon (CZ) potential $u_{\mathrm{CZ}}$
of \cite{CZ2023} is related to $\varphi$ by the explicit dictionary \textbf{(c)}
\texttt{[defect\_semiclassics.pdf]}:
\begin{equation}\label{eq:cz_dict}
  u_{\mathrm{CZ}} = -3\varphi = -6(\ln\tau)_x, \qquad
  u = -\frac{1}{3}(u_{\mathrm{CZ}})_x, \qquad
  \mathcal{L}_{\mathrm{CZ}} = 9\,\mathcal{L}.
\end{equation}

\paragraph{One-soliton and KP-root data.}
A one-soliton of \eqref{eq:good_bq} with amplitude parameter
$\alpha \in (0,\pi/2)$ and direction sign $\epsilon \in \{+1,-1\}$ has \textbf{(a)}
\begin{equation}\label{eq:sol_data}
  k = \sin\alpha, \quad v = \epsilon\cos\alpha, \quad \omega = kv, \quad
  \ell = -\epsilon\cos 3\alpha, \quad M = 2k, \quad \Pi = -Mv, \quad
  P = \frac{2}{3}\epsilon\,k^3\cos\alpha.
\end{equation}
Its two KP roots \textbf{(a,c)}
\begin{equation}\label{eq:kp_roots_sol}
  p_a(\alpha,\epsilon) = \frac{k + v/\sqrt{3}}{2}
  = \frac{1}{\sqrt{3}}\cos\!\left(\alpha - \frac{\epsilon\pi}{6}\right), \qquad
  p_b(\alpha,\epsilon) = \frac{v/\sqrt{3} - k}{2}
  = -\frac{1}{\sqrt{3}}\cos\!\left(\alpha + \frac{\epsilon\pi}{6}\right)
\end{equation}
are two of the three roots $\{q_1,q_2,q_3\}$ of the depressed KP cubic
\begin{equation}\label{eq:kp_cubic}
  p^3 - \frac{1}{4}p = \frac{\ell}{4\sqrt{3}}
  = -\frac{\epsilon}{4\sqrt{3}}\cos 3\alpha, \qquad
  q_1 + q_2 + q_3 = 0, \quad q_1 q_2 + q_2 q_3 + q_3 q_1 = -\frac{1}{4},
\end{equation}
ordered as $p_a > p_b$ so that $k = p_a - p_b > 0$ and $v = \sqrt{3}(p_a + p_b)$.
The three roots of \eqref{eq:kp_cubic} are
\begin{equation}\label{eq:three_roots}
  q_r(\alpha,\epsilon) = \frac{1}{\sqrt{3}}\cos\!\left(\phi_r^{(\epsilon)} + \frac{\pi}{6}\right)
  = -\frac{1}{\sqrt{3}}\sin\!\left(\phi_r^{(\epsilon)} - \frac{\pi}{3}\right),
\end{equation}
and the three root differences are $\sin\phi_r^{(\epsilon)}$ ($r=1,2,3$), where
\begin{equation}\label{eq:phi_r}
  (\phi_1,\phi_2,\phi_3) := \left(\alpha + \frac{2\pi}{3},\; \alpha - \frac{2\pi}{3},\; \alpha\right),
  \qquad
  \phi_r^{(\epsilon)} := \begin{cases}
    \phi_r, & \epsilon = +1,\\
    \pi - \phi_r, & \epsilon = -1,
  \end{cases}
  \qquad (c_1,c_2,c_3) := (1,1,-2).
\end{equation}

\paragraph{The Corrigan--Zambon type-I defect and its KP vertex.}
In CZ variables $u_{\mathrm{CZ}} = -3\varphi$ ($x < 0$) and $v_{\mathrm{CZ}} = -3\psi$
($x > 0$), the type-I integrable defect at $x=0$ is defined by four sewing
conditions (CZ eqs.~(2.4), (3.15)--(3.18) with $\alpha_{\mathrm{CZ}} = 1/(2\sqrt{3})$,
$\beta_{\mathrm{CZ}} = \eta/(2\sqrt{3})$, $\eta \in \{+1,-1\}$, and real defect
parameter $\sigma \in \mathbb{R}$) \textbf{(a)} \cite{CZ2023}. Defining the
signed defect parameter and KP vertex \textbf{(c)} \texttt{[defect\_semiclassics.pdf]}
\begin{equation}\label{eq:s_qD}
  s := \eta\sigma \in \mathbb{R}, \qquad
  q_D := -\frac{s}{2\sqrt{3}} = -\frac{\eta\sigma}{2\sqrt{3}},
\end{equation}
a one-soliton $\tau_L = 1 + E_1$ ($x<0$), $\tau_R = 1 + z_1 E_1$ ($x>0$) with
$E_1 = \ee^{k_1(x - v_1 t - x_0)}$ crosses the defect without reflection and
acquires the left-to-right Hirota transmission factor \textbf{(a,c)} \cite{CZ2023}
\begin{equation}\label{eq:z1_sol}
  z_1(\alpha_1,\epsilon_1;q_D,\eta)
  = \left(\frac{k_1 - v_1/\sqrt{3} - s/\sqrt{3}}{k_1 + v_1/\sqrt{3} + s/\sqrt{3}}\right)^{-\eta}
  = \left(\frac{p_{1,b} - q_D}{p_{1,a} - q_D}\right)^{\eta}.
\end{equation}
For a soliton of direction $\epsilon_1 = \pm 1$ incident from $x \to -\epsilon_1\infty$
and transmitted to $x \to +\epsilon_1\infty$, the directional transmission factor
is $z_t = z_1^{\epsilon_1} = \big((p_{1,b}-q_D)/(p_{1,a}-q_D)\big)^{\epsilon_1\eta}$
and the spatial shift is $\Delta x_D = -k_1^{-1}\ln|z_t|$.

\paragraph{What this write-up resolves (Route E, step 2).}
In Phase~2d (\texttt{defect\_semiclassics.pdf}) it was shown that at tree level
($O(\hb^{-1})$), a defect sits at a single KP vertex $q_D$, a soliton $S = \{p_a,p_b\}$
is a chord between two vertices of a KP triple $\{q_1,q_2,q_3\}$, and absorbing
$S_i = \{q_i,q_{i+1}\}$ shifts the defect vertex $q_i \to q_{i+1}$ while shifting
the tree-level transmission phase of any test soliton $3$ by the two-soliton
phase $\delta_0(3,S_i) + M_{S_i}\Pi_3$ up to a sign. Meanwhile, in Phase~2e
(\texttt{oneloop\_phase.pdf}), we computed the exact one-loop ($O(\hb^0)$)
two-soliton phase directly from the continuum phonon phase shifts and found that
the raw Faddeev--Korepin integral $I_{12}$ and its antisymmetrised partner
$J_{12}$ differ only by a direction-dependent constant $+\frac{\pi}{2}\delta_{\epsilon_1,\epsilon_2}$
from the translational zero-mode jump $\theta_S(0^+) - \theta_S(0^-) = -2\pi$,
whereas both have zero head-on constant and rule out the $\Gamma$-ansatz for
head-on collisions.

Here we carry out \textbf{Option 1 (Route E, step 2)}: we compute the exact
one-loop soliton--defect transmission phase on the CZ defect, solve the defect
vacuum stability spectrum in closed algebraic form, establish exact one-loop
energy--momentum balance in soliton absorption, and test KP-triangle defect
fusion at one loop. We show that one-loop defect fusion holds identically across
every KP absorption edge and closes around every KP triangle \emph{if and only if}
the one-loop two-soliton phase uses the raw Faddeev--Korepin integral $I_{12}$
(without the $+\frac{\pi}{2}\delta_{\epsilon_1,\epsilon_2}$ boundary constant of
$J_{12}$).

% ==============================================================================
\section{Exact time-dependent phonon eigenmodes on $N$-soliton + defect backgrounds}
\label{sec:scattering}
% ==============================================================================

\subsection{Three-soliton Hirota solution on the CZ defect}

Corrigan and Zambon gave explicit Hirota formulas for $1$- and $2$-soliton
solutions on the type-I defect (CZ eqs.~(5.3), (5.6)) \textbf{(a)} \cite{CZ2023}.
By the standard Hirota dressing of B\"acklund/defect transformations, the
$N$-soliton solution on the defect is obtained by replacing each soliton weight
$E_j = \ee^{k_j x - \omega_j t + \delta_j}$ in $\tau_L(E_1,\dots,E_N)$ ($x<0$) by
$z_j E_j$ in $\tau_R(z_1 E_1,\dots,z_N E_N)$ ($x>0$), with the bulk interaction
coefficients $A_{ij} = A(i,j)$ unchanged. For $N=3$,
\begin{equation}\label{eq:tau3_defect}
  \tau_L = 1 + \sum_{j=1}^3 E_j + \sum_{1\le i<j\le 3} A_{ij} E_i E_j
  + A_{12}A_{13}A_{23}\,E_1 E_2 E_3, \qquad
  \tau_R(E_1,E_2,E_3) = \tau_L(z_1 E_1, z_2 E_2, z_3 E_3).
\end{equation}
Substituting $u_{\mathrm{CZ}} = -6(\ln\tau_L)_x$ and $v_{\mathrm{CZ}} = -6(\ln\tau_R)_x$
into the four full nonlinear CZ sewing conditions at $x=0$ (CZ eqs.~(3.15)--(3.18)
in \texttt{verify\_defect\_solutions.py}), we verify at generic parameters that
all four nonlinear sewing conditions hold identically at $35$-digit precision
(maximum residual $1.20\times 10^{-34}$) \textbf{(c)}
\texttt{[oneloop\_defect\_scattering\_and\_stability.py]}.

\subsection{Time-dependent phonon scattering on $1$- and $2$-soliton + defect backgrounds}

To obtain the exact continuum phonon eigenmodes on the defect vacuum ($N=0$),
the $1$-soliton + defect background ($N=1$), and the $2$-soliton + defect
background ($N=2$), we analytically continue the last soliton parameter to a
dispersive plane wave \textbf{(c)}
\texttt{[oneloop\_defect\_scattering\_and\_stability.py]}:
\begin{equation}\label{eq:phonon_continuation}
  k_{\mathrm{ph}} = \ii q, \qquad \omega_{\mathrm{ph}} = \ii\omega(q), \qquad
  \omega(q) := |q|\sqrt{1+q^2}, \qquad
  v_{\mathrm{ph}}(q) = \frac{\omega(q)}{q} = \sgn(q)\sqrt{1+q^2} \qquad (q \in \mathbb{R}\setminus\{0\}),
\end{equation}
and expand $\varphi_{L,R}(x,t) = 2\partial_x(\ln\tau_{L,R}) = \varphi_{L,R}^{(0)}(x,t) + \varepsilon\,\delta\varphi_{L,R}(x,t;q) + O(\varepsilon^2)$
to first order in $E_{\mathrm{ph}} = \varepsilon\,\ee^{\ii(qx-\omega(q)t)}$.

\paragraph{Claim 2.1 (Exact reflectionless phonon eigenmodes on $N$-soliton + defect backgrounds, (c)).}
\texttt{[oneloop\_defect\_scattering\_and\_stability.py]}
\emph{For $N \in \{0,1,2\}$ solitons cross-sewing the CZ defect at $x=0$:
\begin{enumerate}
\item The linearised Hirota field $(\delta\varphi_L(x,t;q),\delta\varphi_R(x,t;q))$
satisfies the bulk linearised good Boussinesq equation
\[
  \big[\partial_t^2 - \partial_x^2 + \partial_x^4 + 6\partial_x(\varphi_{L,R}^{(0)})_x\partial_x\big]\delta\varphi_{L,R}(x,t;q) = 0
\]
for $x<0$ and $x>0$, together with all four linearised CZ sewing conditions at
$x=0$ for every $t \in \mathbb{R}$ (verified at $35$-digit precision with
maximum sewing residual $3.94\times 10^{-34}$ for $N=1$ and $2.10\times 10^{-34}$
for $N=2$ across interacting $t=0.15$ and asymptotic $t=\pm 15$ times).
\item Away from the soliton cores and the defect ($|x| \to \infty$), the
eigenmode is purely oscillatory with \textbf{zero reflection} ($R(q)=0$) and
\textbf{zero asymptotic evanescent waves}:
\begin{equation}\label{eq:asympt_scattering}
  \delta\varphi_L(x,t;q) \underset{x\to-\infty}{\sim} 2\ii q\,\ee^{\ii(qx-\omega(q)t)}, \qquad
  \delta\varphi_R(x,t;q) \underset{x\to+\infty}{\sim} 2\ii q\,T_{\mathrm{tot}}(q)\,\ee^{\ii(qx-\omega(q)t)},
\end{equation}
where the total left-to-right transmission amplitude factors identically for all
$t \in \mathbb{R}$ into the product of the individual soliton and defect
transmission amplitudes:
\begin{equation}\label{eq:T_tot}
  T_{\mathrm{tot}}(q) = \left(\prod_{j=1}^N A\big(j,\mathrm{ph}(q)\big)\right) z_{\mathrm{ph}}(q;q_D,\eta),
  \qquad |T_{\mathrm{tot}}(q)| = 1.
\end{equation}
Here $A(j,\mathrm{ph}(q)) = \ee^{\ii\theta_S(q;\alpha_j,\epsilon_j)}$ is the bulk
soliton--phonon transmission amplitude of \texttt{oneloop\_phase.pdf} and
\begin{equation}\label{eq:z_ph}
  z_{\mathrm{ph}}(q;q_D,\eta)
  = \left(\frac{\sgn(q)\frac{\sqrt{1+q^2}}{\sqrt{3}} + \frac{s}{\sqrt{3}} - \ii q}
               {\sgn(q)\frac{\sqrt{1+q^2}}{\sqrt{3}} + \frac{s}{\sqrt{3}} + \ii q}\right)^{\eta}
  = \left(\frac{\frac{\sgn(q)\sqrt{1+q^2}}{2\sqrt{3}} - q_D - \frac{\ii q}{2}}
               {\frac{\sgn(q)\sqrt{1+q^2}}{2\sqrt{3}} - q_D + \frac{\ii q}{2}}\right)^{\eta}
  =: \ee^{\ii\theta_D(q;s,\eta)}
\end{equation}
is the defect--phonon transmission amplitude \textbf{(c)}
\texttt{[defect\_transmission\_phase.py, oneloop\_defect\_scattering\_and\_stability.py]}.
\end{enumerate}}

\paragraph{Consequence for the vacuum-subtracted one-loop phase.}
On the $1$-soliton + defect background, the total phonon phase shift is the sum
$\theta_{\mathrm{tot}}(q) = \theta_S(q;\alpha,\epsilon) + \theta_D(q;s,\eta)$,
independent of time $t$. Subtracting the defect vacuum ($N=0$, phase shift
$\theta_D(q;s,\eta)$) and the isolated soliton ($N=1$ without defect, phase shift
$\theta_S(q;\alpha,\epsilon)$) leaves only the bilinear cross-interaction
between $\theta_S(q;\alpha,\epsilon)$ and $\theta_D(q;s,\eta)$ in the
Faddeev--Korepin trace formula (Section~\ref{sec:phase}). Likewise, on the
$2$-soliton + defect background,
$\theta_{\mathrm{tot}}(q) = \theta_S(q;\alpha_1,\epsilon_1) + \theta_S(q;\alpha_2,\epsilon_2) + \theta_D(q;s,\eta)$,
so the one-loop $2$-soliton + defect phase is the exact sum of pairwise
phases $\delta_{\mathrm{tot}}^{(1)}(1,2;D) = \delta_{12}^{(1)} + \delta_{1D}^{(1)} + \delta_{2D}^{(1)}$
with zero three-body irreducible one-loop correction \textbf{(c)}.

% ==============================================================================
\section{KP-vertex factorisation and the 2-root dipole decomposition of $\theta_D'(q;s,\eta)$}
\label{sec:dipole}
% ==============================================================================

\subsection{Single-KP-root phonon factor $\mathcal{Z}(q;p)$}

Under the phonon continuation \eqref{eq:phonon_continuation}, the two KP roots
\eqref{eq:kp_roots_sol} become the complex-conjugate pair \textbf{(c)}
\begin{equation}\label{eq:p_ph}
  p_{\mathrm{ph},a}(q) := \frac{\sgn(q)\frac{\sqrt{1+q^2}}{\sqrt{3}} + \ii q}{2}, \qquad
  p_{\mathrm{ph},b}(q) := \frac{\sgn(q)\frac{\sqrt{1+q^2}}{\sqrt{3}} - \ii q}{2} = \overline{p_{\mathrm{ph},a}(q)}.
\end{equation}
For any real or complex KP root $p$, define the elementary vertex--phonon ratio
\textbf{(c)} \texttt{[oneloop\_defect\_scattering\_and\_stability.py]}:
\begin{equation}\label{eq:Z_vertex}
  \mathcal{Z}(q;p) := \frac{p_{\mathrm{ph},b}(q) - p}{p_{\mathrm{ph},a}(q) - p}
  = \frac{\frac{\sgn(q)\sqrt{1+q^2}}{2\sqrt{3}} - p - \frac{\ii q}{2}}
         {\frac{\sgn(q)\sqrt{1+q^2}}{2\sqrt{3}} - p + \frac{\ii q}{2}}.
\end{equation}

\paragraph{Claim 3.1 (KP-vertex/chord factorisation of phonon scattering, (c)).}
\texttt{[oneloop\_defect\_scattering\_and\_stability.py]}
\emph{For every $q \in \mathbb{R}\setminus\{0\}$, every defect vertex $q_D$, and
every soliton $S = \{p_a,p_b\}$, the defect and soliton phonon transmission
amplitudes factor through $\mathcal{Z}(q;p)$ as
\begin{equation}\label{eq:vertex_chord_fact}
  z_{\mathrm{ph}}(q;q_D,\eta) = \mathcal{Z}(q;q_D)^\eta, \qquad
  A\big(\mathrm{sol}(p_a,p_b),\mathrm{ph}(q)\big)
  = \frac{\mathcal{Z}(q;p_b)}{\mathcal{Z}(q;p_a)}
  = z_{\mathrm{ph}}(q;p_b,+1)\,z_{\mathrm{ph}}(q;p_a,-1).
\end{equation}
Consequently, the soliton--phonon phase shift is the exact difference of two
defect--phonon phase shifts at the endpoints $p_b$ and $p_a$ of the soliton
chord (verified to $6.48\times 10^{-36}$):
\begin{equation}\label{eq:theta_S_as_diff_theta_D}
  \theta_S'(q;\alpha,\epsilon) = \theta_D'\big(q; s(p_b), +1\big) - \theta_D'\big(q; s(p_a), +1\big),
  \qquad s(p) := -2\sqrt{3}\,p.
\end{equation}}

\subsection{The 2-root dipole decomposition of $\theta_D'(q;s,\eta)$}

In Phase~2e (\texttt{oneloop\_phase.pdf}, Claim~2.2) we showed that the
soliton--phonon phase-shift derivative $\theta_S'(q;\alpha,\epsilon)$ decomposes
into a sum over the three roots of its KP cubic with weights $(c_1,c_2,c_3) = (1,1,-2)$:
\begin{equation}\label{eq:theta_S_3root}
  \theta_S'(q;\alpha,\epsilon) = \sum_{r=1}^3 c_r\,\Theta'\!\left(q;\phi_r^{(\epsilon)}\right),
  \qquad
  \Theta'(q;\psi) := \frac{\sin\psi}{q^2+\sin^2\psi}
  + \frac{\sgn(q)}{\sqrt{1+q^2}}\frac{\sin\psi\cos\psi}{q^2+\sin^2\psi}.
\end{equation}
Why did a soliton have weights $(1,1,-2)$? Claim~3.2 below reveals the reason:
each single KP vertex $q_D$ is a \emph{2-root dipole} with weights $(-1,+1)$
associated with the two chords connecting $q_D$ to the other two vertices
$q_\pm$ of its KP cubic, and subtracting the two dipoles at $p_b$ and $p_a$
gives $(+1,-1,0) - (0,+1,-1) = (+1,-2,+1)$!

\paragraph{Claim 3.2 (2-root dipole decomposition of the defect phase-shift derivative, (c)).}
\texttt{[oneloop\_defect\_scattering\_and\_stability.py]}
\emph{Given $s = \eta\sigma \in \mathbb{R}$ and $q_D = -s/(2\sqrt{3})$, define
\begin{equation}\label{eq:gamma_psi_def}
  \gamma := \arccos\!\left(\frac{s}{2}\right), \qquad
  \psi_1 := \frac{2\pi}{3} - \gamma, \qquad
  \psi_2 := -\frac{2\pi}{3} - \gamma, \qquad
  (c_{D,1},c_{D,2}) := (-\eta,+\eta),
\end{equation}
where for $|s|\le 2$ ($|q_D|\le 1/\sqrt{3}$) the angle $\gamma \in [0,\pi]$ is
real, while for $|s|>2$ ($\gamma$ complex) $\psi_{1,2}$ are complex conjugates
modulo $\pi$. Then:
\begin{enumerate}
\item The vertex $q_D = -(\cos\gamma)/\sqrt{3}$ belongs to the KP cubic
$4q^3 - q = 4q_D^3 - q_D = -\frac{\cos 3\gamma}{3\sqrt{3}}$, whose other two
roots are $q_\pm = \frac{-q_D \pm \sqrt{1-3q_D^2}}{2} = \frac{-q_D \pm \sin\gamma}{2}$,
and the two root differences from $q_D$ are identically
\begin{equation}\label{eq:sin_psi_roots}
  \sin\psi_1 = \sin\!\left(\gamma + \frac{\pi}{3}\right) = q_+ - q_D, \qquad
  \sin\psi_2 = \sin\!\left(\gamma - \frac{\pi}{3}\right) = q_D - q_-.
\end{equation}
\item For every $q \in \mathbb{R}\setminus\{0\}$ and every $s \in \mathbb{R}$
(both $|s|\le 2$ and $|s|>2$), the defect phase-shift derivative
$\theta_D'(q;s,\eta) := -\ii\partial_q\ln z_{\mathrm{ph}}(q;q_D,\eta)$ is given
by the 2-root dipole sum (verified to $1.88\times 10^{-36}$):
\begin{equation}\label{eq:theta_D_dipole}
  \theta_D'(q;s,\eta)
  = -\eta\,\frac{2\,\sgn(q)\frac{q}{2\sqrt{3}\sqrt{1+q^2}}\frac{q}{2}
    + \left(\sgn(q)\frac{\sqrt{1+q^2}}{2\sqrt{3}} - q_D\right)}
    {\left(\sgn(q)\frac{\sqrt{1+q^2}}{2\sqrt{3}} - q_D\right)^2 + \frac{q^2}{4}}
  = \Rea\sum_{s'=1}^2 c_{D,s'}\,\Theta'(q;\psi_{s'}).
\end{equation}
\end{enumerate}}

\begin{proof}[Derivation]
Rationalising the denominator of \eqref{eq:theta_D_dipole},
\[
  \left(\sgn(q)\frac{\sqrt{1+q^2}}{2\sqrt{3}} - q_D\right)^2 + \frac{q^2}{4}
  = \frac{1 + 4q^2 + 3s^2 + 2s\,\sgn(q)\sqrt{1+q^2}}{12}.
\]
Multiplying numerator and denominator by $(1+4q^2+3s^2) - 2s\,\sgn(q)\sqrt{1+q^2}$
produces
\[
  (1 + 4q^2 + 3s^2)^2 - 4s^2(1 + q^2)
  = 16\left(q^2 + \sin^2\!\left(\gamma+\frac{\pi}{3}\right)\right)
      \left(q^2 + \sin^2\!\left(\gamma-\frac{\pi}{3}\right)\right)
  = 16(q^2 + \sin^2\psi_1)(q^2 + \sin^2\psi_2).
\]
Partial-fraction decomposition in $q^2$ then splits the rational part into
$\frac{-\sin\psi_1}{q^2+\sin^2\psi_1} + \frac{\sin\psi_2}{q^2+\sin^2\psi_2}$ and
the $\sgn(q)/\sqrt{1+q^2}$ part into
$\frac{-\sin\psi_1\cos\psi_1}{q^2+\sin^2\psi_1} + \frac{\sin\psi_2\cos\psi_2}{q^2+\sin^2\psi_2}$,
which is precisely \eqref{eq:theta_D_dipole}.
\end{proof}

% ==============================================================================
\section{Closed-form algebraic solution of the defect vacuum growing-mode spectrum}
\label{sec:stability}
% ==============================================================================

In \texttt{corrigan-zambon/defect\_vacuum\_stability.py}, the linearised
stability of the defect vacuum $\varphi_L^{(0)}=\varphi_R^{(0)}=0$ was studied
numerically by searching for roots of a $4\times 4$ sewing determinant, revealing
a threshold at $s = \eta\sigma = -1$, a second threshold at $s = +1$, and a mode
collision at $s = +2$. Using the KP-vertex structure, we now factor the $4\times 4$
determinant in closed algebraic form \textbf{(c)}
\texttt{[oneloop\_defect\_scattering\_and\_stability.py]}.

\paragraph{Claim 4.1 (Closed-form factorisation of the localised-mode determinant and exact growing-mode eigenvalues, (c)).}
\texttt{[oneloop\_defect\_scattering\_and\_stability.py]}
\emph{Consider a localised eigenmode of frequency $\omega = \ii\gamma_{\mathrm{grow}}$
($\gamma_{\mathrm{grow}} > 0$), which decays as $\ee^{+\kappa_j x}$ for $x<0$ and
$\ee^{-\kappa_j x}$ for $x>0$, where $\kappa_{1,2}^2 = \frac{1 \pm \sqrt{1-4\gamma_{\mathrm{grow}}^2}}{2}$
satisfy $\kappa_1\kappa_2 = \gamma_{\mathrm{grow}}$ and
\begin{equation}\label{eq:S_def}
  S := \kappa_1 + \kappa_2 = \sqrt{1 + 2\gamma_{\mathrm{grow}}} > 1
  \qquad \Longleftrightarrow \qquad
  \gamma_{\mathrm{grow}} = \frac{S^2 - 1}{2}.
\end{equation}
Then:
\begin{enumerate}
\item After dividing out the trivial root $\kappa_1 = \kappa_2$, the $4\times 4$
linearised CZ sewing determinant factors in closed form as (verified to
$3.89\times 10^{-36}$):
\begin{equation}\label{eq:det_factorisation}
  \det M_{\mathrm{red}}(S;\sigma,\eta)
  := \frac{\det M_{4\times 4}(\gamma_{\mathrm{grow}};\sigma,\eta)}{(\kappa_1-\kappa_2)^2}
  = -\frac{S^2(S^2-1)}{24}\,P_-(S;\eta\sigma)\,P_+(S;\eta\sigma),
\end{equation}
where $P_\mp(S;s)$ ($s=\eta\sigma$) are the two quadratic polynomials
\begin{equation}\label{eq:P_mp}
  P_\mp(S;s) := S^2 - \frac{\sqrt{3}\mp 1}{2}\,s\,S + \frac{s^2 - 2 \mp \sqrt{3}}{2}.
\end{equation}
\item The discriminants of $P_\mp(S;s)$ are perfect squares times $(4-s^2)$:
\begin{equation}\label{eq:disc_P}
  \Delta_\mp(s) = \left(\frac{\sqrt{3}\pm 1}{2}\right)^2 (4 - s^2),
\end{equation}
yielding the four roots
\begin{equation}\label{eq:S_roots}
  S_{-,\pm}(s) = \frac{(\sqrt{3}-1)s \pm (\sqrt{3}+1)\sqrt{4-s^2}}{4}, \qquad
  S_{+,\pm}(s) = \frac{(\sqrt{3}+1)s \pm (\sqrt{3}-1)\sqrt{4-s^2}}{4}.
\end{equation}
\item Substituting $S_{-,-}(s)$ and $S_{+,+}(s)$ into $\gamma_{\mathrm{grow}} = (S^2-1)/2$
gives the \textbf{exact closed-form localised defect eigenvalues} (verified
against numerical root-finding of $\det M_{4\times 4}=0$ to $1.62\times 10^{-35}$):
\begin{equation}\label{eq:omega_closed}
  \boxed{\;
  \omega_1(s) = \ii\,\frac{s\sqrt{4-s^2} - \sqrt{3}(2-s^2)}{8}, \qquad
  \omega_2(s) = \ii\,\frac{s\sqrt{4-s^2} + \sqrt{3}(2-s^2)}{8}
  \qquad (s = \eta\sigma).
  \;}
\end{equation}
\end{enumerate}}

\paragraph{Complete analytic stability phase diagram and Levinson mode count \textbf{(c)}.}
\texttt{[oneloop\_defect\_scattering\_and\_stability.py]}
Formula \eqref{eq:det_factorisation}--\eqref{eq:omega_closed} accounts analytically
for every regime of the CZ defect vacuum:
\begin{itemize}
\item \textbf{Strictly stable half-line $s \le -1$ ($q_D \ge \frac{1}{2\sqrt{3}}$):}
Neither $S_{-,-}(s)$ nor $S_{+,+}(s)$ exceeds $1$, so there are \textbf{zero}
localised growing or bound modes. For $s \in [-2,-1)$, $\gamma \in (2\pi/3,\pi]$,
so $\psi_1 = 2\pi/3-\gamma \in [-\pi/3,0)$ has $\sin\psi_1 < 0$ while
$\psi_2 = -2\pi/3-\gamma \in [-5\pi/3,-4\pi/3)$ has $\sin\psi_2 > 0$. Since
$\int_{-\infty}^\infty \Theta'(q;\psi)\,\dd q = 2\pi\sgn(\sin\psi)$ is an even
function of $q$ while the $\sgn(q)/\sqrt{1+q^2}$ part is odd and integrates to
$\pm\frac{2\pi}{3}$ on $q>0$ and $q<0$ separately, the total Levinson phase
winding across $q \in (-\infty,\infty)$ is
\begin{equation}\label{eq:levinson_defect}
  \Delta\theta_D(s,\eta=1) := \int_{-\infty}^\infty \theta_D'(q;s,1)\,\dd q
  = \pi\big(-\sgn(\sin\psi_1) + \sgn(\sin\psi_2)\big)
  = \begin{cases}
    +2\pi, & s < -1,\\
    0, & -1 < s < +1,\\
    -2\pi, & s > +1,
  \end{cases}
\end{equation}
decreasing by $2\pi$ each time a growing mode bifurcates out of $\omega=0$.
\item \textbf{First threshold $s = -1$ ($q_D = +\frac{1}{2\sqrt{3}}$, $\gamma = 2\pi/3$):}
Here $\psi_1 = 0$, so $\sin\psi_1$ changes sign from negative to positive,
$P_-(1;-1) = 0$ ($S_{-,-}(-1) = 1$, $\omega_1(-1) = 0$), and the first growing
mode $\omega_1(s) = \ii\gamma_1(s)$ bifurcates out of the $\omega=0$ continuum
edge for $s > -1$, accompanied by a jump of $-2\pi$ in $\Delta\theta_D$ (from
$+2\pi$ to $0$).
\item \textbf{Second threshold $s = +1$ ($q_D = -\frac{1}{2\sqrt{3}}$, $\gamma = \pi/3$):}
Here $\psi_2 = -\pi$, so $\sin\psi_2$ changes sign from positive to negative,
$P_+(1;+1) = 0$ ($S_{+,+}(+1) = 1$, $\omega_2(+1) = 0$), and the second growing
mode $\omega_2(s) = \ii\gamma_2(s)$ bifurcates out of $\omega=0$ for $s > +1$,
accompanied by a second jump of $-2\pi$ in $\Delta\theta_D$ (from $0$ to $-2\pi$).
\item \textbf{Collision threshold $s = +2$ ($q_D = -1/\sqrt{3}$, $\gamma = 0$) and oscillatory regime $s > 2$:}
At $s = +2$, the discriminant $4-s^2$ vanishes and the two pure-imaginary
eigenvalues collide at $\omega_1(2) = \omega_2(2) = \ii\sqrt{3}/4$ ($S = (\sqrt{3}+1)/2$).
For $s > 2$, $\sqrt{4-s^2} = \pm\ii\sqrt{s^2-4}$ and \eqref{eq:omega_closed}
becomes the complex-conjugate oscillatory unstable pair
\[
  \omega_\pm(s) = \frac{\pm s\sqrt{s^2-4} + \ii\sqrt{3}(s^2-2)}{8}.
\]
\end{itemize}

\paragraph{Why localised defect modes cancel in the soliton--defect transmission phase \textbf{(c)}.}
In the strictly stable regime $s \le -1$, no localised defect modes exist at all.
Even when $s > -1$, a soliton crosses the defect at $x=0$ in a finite time
interval and leaves the asymptotic defect state at $t \to +\infty$ identical to
that at $t \to -\infty$ (unless $q_D = p_{a,b}$, which is an inelastic absorption
threshold). Because the one-loop soliton--defect transmission phase is the
vacuum-subtracted difference between the soliton+defect sector and the
isolated-soliton plus defect-vacuum sectors, any localised mode attached to the
static defect has the same eigenfrequency before and after the collision and
cancels between the soliton+defect and defect-vacuum functional determinants.

% ==============================================================================
\section{Closed-form $3\times 2$ dilogarithm formula for the one-loop soliton--defect phase and $(E,P)$ absorption balance}
\label{sec:phase}
% ==============================================================================

\subsection{Equality of raw and antisymmetrised Faddeev--Korepin integrals: $I_{SD} = J_{SD}^{\mathrm{phys}}$}

In Phase~2e (\texttt{oneloop\_phase.pdf}), the raw Faddeev--Korepin integral
$I_{12}$ and the antisymmetrised integral $J_{12}$ for two solitons differed by
$J_{12} - I_{12} = \frac{\pi}{2}\delta_{\epsilon_1,\epsilon_2}$ because the
soliton has two discrete zero modes (translational and $\alpha$-derivative), giving
$\theta_S(0^+) - \theta_S(0^-) = -2\pi$ and a non-vanishing integration-by-parts
boundary term at $q = 0^\pm$.

For a soliton scattering on a \emph{static} defect at $x=0$, the defect breaks
translation invariance explicitly and has \textbf{no translational zero mode}.
Indeed, from \eqref{eq:z_ph}, as $q \to 0^\pm$,
\begin{equation}\label{eq:z_ph_zero}
  z_{\mathrm{ph}}(0^\pm;q_D,\eta)
  = \left(\frac{\frac{\pm 1}{2\sqrt{3}} - q_D - \ii\,0^+}{\frac{\pm 1}{2\sqrt{3}} - q_D + \ii\,0^+}\right)^{\eta}
  = 1 \qquad \left(q_D \ne \pm\frac{1}{2\sqrt{3}}\right),
\end{equation}
so the physical defect phase shift $\theta_D^{\mathrm{phys}}(q;s,\eta) := -\ii\ln z_{\mathrm{ph}}(q;q_D,\eta)$
defined continuously from $q=0^\pm$ vanishes at both long-wavelength limits:
\begin{equation}\label{eq:theta_D_zero}
  \theta_D^{\mathrm{phys}}(0^+;s,\eta) = \theta_D^{\mathrm{phys}}(0^-;s,\eta) = 0, \qquad
  \theta_S(\pm\infty;\alpha,\epsilon) = 0.
\end{equation}
Integrating by parts on $(-\infty,0)$ and $(0,\infty)$ therefore yields zero
boundary terms at both $q=0^\pm$ and $q=\pm\infty$!

\paragraph{Claim 5.1 (Equality $I_{SD} = J_{SD}^{\mathrm{phys}}$ and closed-form $3\times 2$ dilogarithm formula, (c)).}
\texttt{[oneloop\_defect\_phase\_and\_fusion.py]}
\emph{For every soliton $(\alpha,\epsilon) \in (0,\pi/2)\times\{+1,-1\}$ and every
defect $(s,\eta) \in (\mathbb{R}\setminus\{\pm 1\})\times\{+1,-1\}$, the raw and
antisymmetrised Faddeev--Korepin one-loop soliton--defect continuum integrals
coincide:
\begin{align}\label{eq:ISD_eq_JSD}
  I_{SD}(\alpha,\epsilon;s,\eta)
  &:= \frac{1}{4\pi}\int_{-\infty}^\infty \theta_S(q;\alpha,\epsilon)\,\theta_D'(q;s,\eta)\,\dd q \nonumber\\
  &= \frac{1}{8\pi}\int_{-\infty}^\infty \Big[\theta_S(q;\alpha,\epsilon)\,\theta_D'(q;s,\eta)
    - \theta_D^{\mathrm{phys}}(q;s,\eta)\,\theta_S'(q;\alpha,\epsilon)\Big]\dd q
  =: J_{SD}^{\mathrm{phys}}(\alpha,\epsilon;s,\eta),
\end{align}
and evaluate in closed analytical form as the $3\times 2$ dilogarithm sum
(verified against numerical quadrature across $|s|<2$ and $|s|>2$ to
$4.51\times 10^{-36}$, and against the Faddeev--Korepin $\alpha$-gradient
$\partial_\alpha I_{SD} = \frac{1}{4\pi}\int (\partial_\alpha\theta_S)\theta_D'\,\dd q$
to $9.40\times 10^{-38}$):
\begin{equation}\label{eq:ISD_closed}
  \boxed{\;
  I_{SD}(\alpha,\epsilon;s,\eta) = J_{SD}^{\mathrm{phys}}(\alpha,\epsilon;s,\eta)
  = \frac{\epsilon}{2\pi}\sum_{r=1}^3\sum_{s'=1}^2 c_r\,c_{D,s'}\,
    \Rea L\!\left(\tan\frac{\phi_r}{2}\left(\tan\frac{\psi_{s'}}{2}\right)^{-\epsilon}\right),
  \;}
\end{equation}
where $(\phi_1,\phi_2,\phi_3) = (\alpha+2\pi/3,\alpha-2\pi/3,\alpha)$,
$(c_1,c_2,c_3) = (1,1,-2)$, $(\psi_1,\psi_2) = (2\pi/3-\gamma,-2\pi/3-\gamma)$
with $\gamma = \arccos(s/2)$, $(c_{D,1},c_{D,2}) = (-\eta,+\eta)$, and
\begin{equation}\label{eq:L_def}
  L(z) = \int_0^z \frac{2\ln t}{1-t^2}\,\dd t
  = -\left[2\chi_2(z) + \ln z\,\ln\frac{1-z}{1+z}\right] \quad (\Rea z > 0), \qquad
  L(-z) = -L(z),
\end{equation}
with $\chi_2(z) = \frac{1}{2}\big(\Li_2(z)-\Li_2(-z)\big)$ the Legendre chi
function. Moreover, $I_{SD}$ obeys the exact parity/crossing and defect-orientation
identities (verified to $2.26\times 10^{-36}$):
\begin{equation}\label{eq:ISD_symmetries}
  I_{SD}(\alpha,-1;s,\eta) = I_{SD}(\alpha,+1;-s,\eta), \qquad
  I_{SD}(\alpha,\epsilon;s,-\eta) = -I_{SD}(\alpha,\epsilon;s,\eta).
\end{equation}}

\subsection{One-loop defect vacuum shifts $(\Delta E_D^{\mathrm{ren}}, \Delta P_D)$ and exact absorption balance}

Because the defect phase-shift derivative $\theta_D'(q;s,\eta)$ is a 2-root
dipole \eqref{eq:theta_D_dipole} over the universal kernel $\Theta'(q;\psi)$
with $\sum_{s'=1}^2 c_{D,s'}\sin\psi_{s'} \ne 0$, the raw zero-point energy shift
of the defect vacuum has a UV divergence proportional to $\sum_{s'} c_{D,s'}\sin\psi_{s'}$
that is cancelled by the exact same bulk cubic tadpole counterterm as the
soliton vacuum energy in \texttt{oneloop\_phase.pdf} (since the defect sewing
potential is linear in $u_{\mathrm{CZ}}$ and has zero bulk tadpole of its own).
Meanwhile, the defect vacuum zero-point momentum shift $\Delta P_D$ is UV
convergent without any counterterm because the odd part of $\theta_D'(q)$ decays
as $O(q^{-3})$.

\paragraph{Claim 5.2 (Closed-form defect vacuum shifts and exact one-loop $(E,P)$ balance in soliton absorption, (c)).}
\texttt{[oneloop\_defect\_phase\_and\_fusion.py]}
\emph{For any defect vertex $q_D = -s/(2\sqrt{3})$ and orientation $\eta \in \{+1,-1\}$:
\begin{enumerate}
\item The one-loop zero-point momentum $\Delta P_D(q_D,\eta) = \int_{-\infty}^\infty \frac{\dd q}{4\pi}\,q\,\theta_D'(q;s,\eta)$
and normal-ordered zero-point energy $\Delta E_D^{\mathrm{ren}}(q_D,\eta)$ are
given in closed form for all $s \in \mathbb{R}$ by (verified against numerical
quadrature to $7.52\times 10^{-37}$):
\begin{align}
  \Delta P_D(q_D,\eta) &= -\frac{1}{2\pi}\sum_{s'=1}^2 c_{D,s'}\,\Rea\!\left[
    \frac{1}{2}\sin\psi_{s'}\ln\frac{1-\cos\psi_{s'}}{1+\cos\psi_{s'}}\right], \label{eq:Delta_PD}\\
  \Delta E_D^{\mathrm{ren}}(q_D,\eta) &= \frac{1}{2\pi}\sum_{s'=1}^2 c_{D,s'}\,\Rea\!\left[
    \sin\psi_{s'}\left(1 - \frac{1}{2}\cos\psi_{s'}\ln\frac{1-\cos\psi_{s'}}{1+\cos\psi_{s'}}\right)\right]. \label{eq:Delta_ED}
\end{align}
\item Let $\{q_1,q_2,q_3\}$ be any KP triple ($|q_i| < 1/\sqrt{3}$), and let
$S_i = \{q_i,q_{i+1}\}$ be the soliton connecting vertex $q_i$ to vertex
$q_{i+1}$ ($q_4 \equiv q_1$). Define the intrinsic absorption sign
\begin{equation}\label{eq:rho_i_def}
  \rho_i := \eta\,\sgn(q_i - q_{i+1}) \in \{+1,-1\}.
\end{equation}
Then across every soliton absorption/emission step $D(q_i) + S_i \leftrightarrow D(q_{i+1})$,
the one-loop soliton energy and momentum shifts
$(\Delta E_1^{\mathrm{ren}}(S_i),\Delta P_1(S_i))$ equal the jump in the
defect vacuum shifts (verified across $18$ absorption edges to
$3.20\times 10^{-36}$):
\begin{equation}\label{eq:ep_balance}
  \boxed{\;
  \Delta E_D^{\mathrm{ren}}(q_{i+1},\eta) - \Delta E_D^{\mathrm{ren}}(q_i,\eta)
  = \rho_i\,\Delta E_1^{\mathrm{ren}}(S_i), \qquad
  \Delta P_D(q_{i+1},\eta) - \Delta P_D(q_i,\eta)
  = \rho_i\,\Delta P_1(S_i).
  \;}
\end{equation}
\end{enumerate}}

% ==============================================================================
\section{One-loop defect fusion, KP-triangle closure, and implications for the quantum S-matrix}
\label{sec:fusion}
% ==============================================================================

\subsection{Total one-loop soliton--defect phase at fixed physical momentum}

Including the shift $\alpha \mapsto \alpha + \hb\,\delta\alpha_1$ required to
hold the physical soliton momentum $P_{\mathrm{phys}} = P(\alpha,\epsilon) + \hb\,\Delta P_1(\alpha,\epsilon)$
fixed at $O(\hb)$ (as in \texttt{oneloop\_phase.pdf}, Section~3.3), where
\begin{equation}\label{eq:dPdalpha_defect}
  \frac{\partial \delta_D^{\mathrm{trans}}(\alpha,\epsilon;s,\eta)}{\partial \alpha}
  = \Delta x_D(\alpha,\epsilon;s,\eta)\,\frac{\dd P(\alpha,\epsilon)}{\dd\alpha}, \qquad
  \Delta x_D(\alpha,\epsilon;s,\eta) = -\frac{\epsilon\eta}{k}\ln\left|\frac{p_b - q_D}{p_a - q_D}\right|,
\end{equation}
 gives the fixed-physical-momentum renormalisation shift
\begin{equation}\label{eq:delta_renorm_SD}
  \delta_{D,\mathrm{renorm}}^{(q_D)}(\alpha,\epsilon;\eta)
  := -\frac{\partial\delta_D^{\mathrm{trans}}}{\partial P}\,\Delta P_1(\alpha,\epsilon)
  = -\Delta x_D(\alpha,\epsilon;s,\eta)\,\Delta P_1(\alpha,\epsilon),
\end{equation}
where we follow the sign convention of \texttt{oneloop\_defect\_phase\_and\_fusion.py}
in tracking $\Delta x_D\,\Delta P_1$.

\subsection{Exact one-loop defect fusion and KP-triangle closure}

We now test whether the tree-level defect fusion relation
\cite{CZ2023} (\texttt{defect\_semiclassics.pdf}, Claim~5) survives at one loop.

\paragraph{Claim 6.1 (Exact one-loop defect fusion and KP-triangle coboundary closure, (c)).}
\texttt{[oneloop\_defect\_phase\_and\_fusion.py]}
\emph{Let $\{q_1,q_2,q_3\}$ be any KP triple, $S_i = \{q_i,q_{i+1}\}$ ($i=1,2,3$)
the three solitons on the edges of the KP triangle, $\rho_i = \eta\sgn(q_i-q_{i+1}) \in \{+1,-1\}$
the intrinsic absorption sign \eqref{eq:rho_i_def}, and $3 = (\alpha_3,\epsilon_3)$
an arbitrary test soliton. Then:
\begin{enumerate}
\item At tree level, the directional transmission phase $\delta_D^{\mathrm{trans}}(3;q_D,\eta)$
obeys
\begin{equation}\label{eq:tree_fusion}
  \delta_D^{\mathrm{trans}}(3;q_{i+1},\eta) - \delta_D^{\mathrm{trans}}(3;q_i,\eta)
  = \rho_{\mathrm{tree},i}\Big(\delta_0(3,S_i) + M_{S_i}\Pi_3\Big), \qquad
  \rho_{\mathrm{tree},i} = \epsilon_3\rho_i \in \{+1,-1\},
\end{equation}
where the factor $\epsilon_3$ reflects the directional transmission exponent
$z_t = z_3^{\epsilon_3}$ for a right-moving ($\epsilon_3=+1$, $L\to R$) versus
left-moving ($\epsilon_3=-1$, $R\to L$) test soliton.
\item At one loop, both the Faddeev--Korepin continuum phase $I_{SD}(3;q_D,\eta)$
and the momentum-renormalisation shift $\Delta x_D(3;q_D,\eta)\,\Delta P_1(\alpha_3,\epsilon_3)$
satisfy \textbf{exact one-loop defect fusion} on every edge $q_i \to q_{i+1}$
(verified across $12$ KP triangles $= 36$ edges to $4.51\times 10^{-36}$ for
$I_{SD}$ and $1.67\times 10^{-33}$ for the momentum-renormalisation term):
\begin{equation}\label{eq:oneloop_fusion}
  \boxed{\;
  I_{SD}(3;q_{i+1},\eta) - I_{SD}(3;q_i,\eta) = \rho_i\,I_{SS}(3,S_i), \qquad
  \Delta x_D(3;q_{i+1},\eta) - \Delta x_D(3;q_i,\eta) = \epsilon_3\rho_i\,\Delta X_3(3,S_i),
  \;}
\end{equation}
where $I_{SS}(3,S_i) = I_{12}(\alpha_3,\epsilon_3;\alpha_{S_i},\epsilon_{S_i})$
is the raw Faddeev--Korepin one-loop two-soliton phase of \texttt{oneloop\_phase.pdf}
and $\Delta X_3(3,S_i) = -k_3^{-1}\ln|A_{12}(3,S_i)|$. Equivalently, for the
directional transmission integral $I_{SD}^{\mathrm{trans}}(3;q_D,\eta) := \epsilon_3\,I_{SD}(3;q_D,\eta)$,
both terms in \eqref{eq:oneloop_fusion} carry the exact same sign
$\rho_{\mathrm{tree},i} = \epsilon_3\rho_i$ as tree level \eqref{eq:tree_fusion}.
\item Summing \eqref{eq:oneloop_fusion} around any closed KP triangle
$q_1 \to q_2 \to q_3 \to q_1$, the raw Faddeev--Korepin one-loop two-soliton
phase $I_{SS}$ and the momentum-renormalisation shift close \textbf{identically
to zero} (verified across all $12$ triangles to $1.50\times 10^{-35}$ and
$7.64\times 10^{-34}$):
\begin{equation}\label{eq:triangle_closure}
  \sum_{i=1}^3 \rho_i\,I_{SS}(3,S_i) = 0, \qquad
  \sum_{i=1}^3 \rho_{\mathrm{tree},i}\,\Delta X_3(3,S_i)\,\Delta P_1(\alpha_3,\epsilon_3) = 0.
\end{equation}
\item By contrast, if one replaces $I_{SS}(3,S_i)$ by the antisymmetrised
integral $J_{SS}(3,S_i) = I_{SS}(3,S_i) + \frac{\pi}{2}\delta_{\epsilon_3,\epsilon_{S_i}}$
(or by any candidate with a direction-dependent $\pm\frac{\pi}{2}$ constant, such
as the $\Gamma$-ansatz), the closed KP triangle acquires a \textbf{non-zero
coboundary obstruction} (verified to $1.35\times 10^{-35}$):
\begin{equation}\label{eq:triangle_obstruction}
  \sum_{i=1}^3 \rho_i\,J_{SS}(3,S_i)
  = \frac{\pi}{2}\sum_{i=1}^3 \rho_i\,\delta_{\epsilon_3,\epsilon_{S_i}}
  \in \left\{0,\; +\frac{\pi}{2},\; -\frac{\pi}{2}\right\}.
\end{equation}
\end{enumerate}}

\begin{proof}[Algebraic proof of Claim 6.1(2)--(4)]
By Claim~3.1, the soliton--phonon phase-shift derivative for $S_i = \{p_{S_i,a},p_{S_i,b}\}$
($p_{S_i,a} = \max(q_i,q_{i+1})$, $p_{S_i,b} = \min(q_i,q_{i+1})$) is the exact
coboundary $\theta_S'(q;S_i) = \theta_D'(q;s(p_{S_i,b}),+1) - \theta_D'(q;s(p_{S_i,a}),+1)$.
Since $\theta_D'(q;s,\eta) = \eta\,\theta_D'(q;s,+1)$ and
$(q_{i+1},q_i) = (p_{S_i,b},p_{S_i,a})$ when $q_{i+1} < q_i$ ($\rho_i = +\eta$)
and $(p_{S_i,a},p_{S_i,b})$ when $q_{i+1} > q_i$ ($\rho_i = -\eta$), we have
pointwise for every $q \in \mathbb{R}\setminus\{0\}$:
\begin{equation}\label{eq:pointwise_fusion}
  \theta_D'(q;s(q_{i+1}),\eta) - \theta_D'(q;s(q_i),\eta) = \rho_i\,\theta_S'(q;S_i).
\end{equation}
Multiplying \eqref{eq:pointwise_fusion} by $\frac{1}{4\pi}\theta_S(q;\alpha_3,\epsilon_3)$
and integrating over $q \in \mathbb{R}$ immediately gives
$I_{SD}(3;q_{i+1},\eta) - I_{SD}(3;q_i,\eta) = \rho_i\,I_{SS}(3,S_i)$. Summing
over $i=1,2,3$ telescopes the left-hand side to zero, proving
$\sum_{i=1}^3\rho_i\,I_{SS}(3,S_i) = 0$. By contrast, the extra boundary term
$\frac{\pi}{2}\delta_{\epsilon_3,\epsilon_{S_i}}$ in $J_{SS}(3,S_i)$ comes from
integrating $\theta_S(q;S_i)\theta_S'(q;3)$ by parts across the soliton zero-mode
jump $\theta_S(0^+)-\theta_S(0^-)=-2\pi$, which has no counterpart in
$\theta_D^{\mathrm{phys}}(0^\pm)=0$ and is \emph{not} a coboundary on the KP
vertices $\{q_1,q_2,q_3\}$.
\end{proof}

\paragraph{What Claim 6.1 implies for the quantum two-soliton S-matrix.}
In Phase~2e (\texttt{oneloop\_phase.pdf}), we noted a $O(\hb^0)$ convention choice
for same-direction collisions ($\epsilon_1 = \epsilon_2$) between the raw
Faddeev--Korepin integral $I_{12}$ (which satisfies $I_{12}(\alpha,\alpha) = 0$,
corresponding to bosonic normalisation $S_{\mathrm{same}}(\alpha,\alpha) = +1$)
and the antisymmetrised integral $J_{12} = I_{12} + \frac{\pi}{2}\delta_{\epsilon_1,\epsilon_2}$
(which approaches $+\frac{\pi}{2}$ as $\alpha_1 \to \alpha_2$, corresponding to
fermionic/exclusion normalisation $S_{\mathrm{same}} \to \ii$). Claim~6.1
provides a decisive algebraic consistency selection rule from integrable defect
fusion:
\begin{quote}
\textbf{Defect-fusion selection rule \textbf{(c)}:} \emph{Only the raw
Faddeev--Korepin one-loop phase $I_{12}$ (together with the momentum-renormalisation
term $-\Delta X_{12}\Delta P$) is an exact coboundary on the space of KP vertices,
$I_{12}(3,S) = I_{SD}(3;p_b,+1) - I_{SD}(3;p_a,+1)$, and therefore satisfies
one-loop defect fusion \eqref{eq:oneloop_fusion} and closes around every KP
triangle \eqref{eq:triangle_closure}.}
\end{quote}

\subsection{Small-amplitude limit, absorption lines, and stability/cusp jumps}

We record the limiting regimes of $I_{SD}(\alpha,\epsilon;s,\eta)$ and
$\delta_{D,\mathrm{renorm}}^{(q_D)}(\alpha,\epsilon;\eta)$ \textbf{(c)}
\texttt{[oneloop\_defect\_phase\_and\_fusion.py]}:
\begin{enumerate}
\item \textbf{Small-amplitude limit $\alpha \to 0$:}
Using $L(w\,\alpha/2) = -w\alpha(\ln(\alpha/2)+\ln|w|-1) + O(\alpha^3\ln\alpha)$
for $\phi_3 = \alpha$ and the Taylor expansion of $L(\tan(\phi_{1,2}/2)w)$ around
$\phi_{1,2} = \pm 2\pi/3$, $I_{SD}(\alpha,\epsilon;s,\eta)$ vanishes as
$O(\alpha\ln\alpha) + O(\alpha)$:
\begin{equation}\label{eq:ISD_small_alpha}
  I_{SD}(\alpha,\epsilon;s,\eta)
  = \frac{\epsilon}{2\pi}\sum_{s'=1}^2 c_{D,s'}\,\Rea\!\Big[
    -2w_{s'}\alpha\big(\ln(\alpha/2) + \ln|w_{s'}| - 1\big)
    + 4w_{s'}\alpha\,L'\big(\sqrt{3}\,w_{s'}\big)\Big]
  + O(\alpha^3\ln\alpha),
\end{equation}
where $w_{s'} := (\tan(\psi_{s'}/2))^{-\epsilon}$ and $L'(z) = \frac{2\ln|z|}{1-z^2}$
(verified at $\alpha = 10^{-2}, 10^{-4}, 10^{-6}$ with error scaling as
$1.57\times 10^{-7}, 3.45\times 10^{-13}, 5.33\times 10^{-19}$).
\item \textbf{Soliton absorption lines $q_D \to p_a(\alpha,\epsilon)$ or $p_b(\alpha,\epsilon)$:}
At an absorption threshold where the soliton chord shares a vertex with the
defect ($z_1 \to 0$ or $\infty$), none of the arguments
$\tan(\phi_r/2)(\tan(\psi_{s'}/2))^{-\epsilon}$ equals $\pm 1$ (since $\phi_r$
encodes the root differences of the soliton triple while $\psi_{s'}$ encodes the
root differences of the defect triple, which only coincide if the two triples are
identical). Consequently, the continuum integral $I_{SD}(\alpha,\epsilon;s,\eta)$
remains \textbf{smooth and finite} across the absorption line (e.g.\ at $s=0.7$,
$\epsilon=-1$, $\alpha_{\mathrm{abs}} \approx 0.659088$,
$I_{SD}(\alpha_{\mathrm{abs}},-1;0.7,1) = 0.7225885601$), whereas the
fixed-physical-momentum shift $\delta_{D,\mathrm{renorm}}^{(q_D)} \propto \ln|z_t|\,\Delta P_1$
diverges logarithmically as $\pm k^{-1}\Delta P_1(\alpha_{\mathrm{abs}},\epsilon)\ln|\alpha-\alpha_{\mathrm{abs}}|$.
\item \textbf{Jumps across the defect stability threshold $s = -\epsilon$ and the soliton cusp $\alpha = \pi/3$:}
When a parameter crosses a locus where $\tan(\phi_r/2)(\tan(\psi_{s'}/2))^{-\epsilon} = \pm 1$,
$\Rea L(z)$ has a finite slope while the integral representation \eqref{eq:ISD_eq_JSD}
jumps by $\pm\pi$ as a pole $q = \pm\sin\psi_{s'}$ crosses $q=0$ or merges with
a soliton pole:
\begin{itemize}
\item Across the defect stability threshold $s = -\epsilon$ ($\gamma = 2\pi/3$ for
$\epsilon=+1$, where $\psi_1=0$, or $\gamma = \pi/3$ for $\epsilon=-1$, where
$\psi_2 = -\pi$), $\tan(\psi_{s'}/2)^{-\epsilon}$ changes sign through $0$ or $\infty$,
and $I_{SD}(\alpha,\epsilon;s,\eta)$ jumps by $-\epsilon\eta\pi$ (verified at
$s = -1\pm 10^{-10}$ for $\epsilon=+1$: jump $-\pi$; and at $s = +1\pm 10^{-10}$
for $\epsilon=-1$: jump $+\pi$).
\item Across the soliton cusp $\alpha = \pi/3$ ($\phi_1 = \pi$, $\tan(\phi_1/2) \to \pm\infty$),
$I_{SD}(\alpha,\epsilon;s,\eta)$ jumps by $+\epsilon\eta\pi$ for $|s|>1$ and is
continuous ($0$ jump) for $|s|<1$ (verified at $s=-1.5$: jump $+\pi$; at $s=0.5$:
jump $0$).
\end{itemize}
\end{enumerate}

% ==============================================================================
\section{Summary of claims, verification scripts, and failures recorded}
\label{sec:summary}
% ==============================================================================

\begin{table}[htbp]
\centering
\footnotesize
\begin{tabular}{>{\raggedright\arraybackslash}p{4.3cm} c >{\raggedright\arraybackslash}p{5.6cm} >{\raggedright\arraybackslash}p{3.8cm}}
\toprule
\textbf{Claim / Formula} & \textbf{Label} & \textbf{Verification script} & \textbf{Precision / Result} \\
\midrule
3-soliton Hirota solution \eqref{eq:tau3_defect} on the CZ defect satisfies all 4 nonlinear sewing conditions at $x=0$ &
\textbf{(c)} & \texttt{oneloop\_defect\_scattering\_and\_stability.py} &
Max residual $1.20\times 10^{-34}$ \\
\addlinespace
Exact time-dependent phonon eigenmodes on $1$- and $2$-soliton + defect backgrounds (Claim~2.1): $R(q)=0$ and factorised $T_{\mathrm{tot}}(q)$ \eqref{eq:T_tot} for all $t \in \mathbb{R}$ &
\textbf{(c)} & \texttt{oneloop\_defect\_scattering\_and\_stability.py} &
Max linearised sewing residual $3.94\times 10^{-34}$ ($N=1$), $2.10\times 10^{-34}$ ($N=2$) \\
\addlinespace
KP-vertex/chord factorisation \eqref{eq:vertex_chord_fact} and 2-root $(-\eta,+\eta)$ dipole decomposition of $\theta_D'(q;s,\eta)$ \eqref{eq:theta_D_dipole} (Claims~3.1--3.2) &
\textbf{(c)} & \texttt{oneloop\_defect\_scattering\_and\_stability.py} &
Max residual $6.48\times 10^{-36}$ across $|s|<2$ and $|s|>2$ \\
\addlinespace
Closed-form factorisation of the $4\times 4$ localised-mode determinant \eqref{eq:det_factorisation}, exact growing-mode eigenvalues \eqref{eq:omega_closed}, and Levinson jumps \eqref{eq:levinson_defect} (Claim~4.1) &
\textbf{(c)} & \texttt{oneloop\_defect\_scattering\_and\_stability.py} &
Det.\ residual $3.89\times 10^{-36}$; root residual $1.62\times 10^{-35}$; Levinson residual $6.12\times 10^{-34}$ \\
\addlinespace
Equality $I_{SD}=J_{SD}^{\mathrm{phys}}$ \eqref{eq:ISD_eq_JSD}, closed-form $3\times 2$ dilogarithm formula \eqref{eq:ISD_closed}, crossing symmetry \eqref{eq:ISD_symmetries}, and FK $\alpha$-gradient (Claim~5.1) &
\textbf{(c)} & \texttt{oneloop\_defect\_phase\_and\_fusion.py} &
Quadrature match $4.51\times 10^{-36}$; crossing $2.26\times 10^{-36}$; FK gradient $9.40\times 10^{-38}$ \\
\addlinespace
Closed-form defect vacuum shifts $(\Delta E_D^{\mathrm{ren}},\Delta P_D)$ \eqref{eq:Delta_PD}--\eqref{eq:Delta_ED} and exact one-loop $(E,P)$ balance in soliton absorption \eqref{eq:ep_balance} (Claim~5.2) &
\textbf{(c)} & \texttt{oneloop\_defect\_phase\_and\_fusion.py} &
Quadrature match $7.52\times 10^{-37}$; balance across $18$ absorption edges $3.20\times 10^{-36}$ \\
\addlinespace
Exact one-loop defect fusion \eqref{eq:oneloop_fusion}, KP-triangle closure for $I_{SS}$ \eqref{eq:triangle_closure}, and coboundary obstruction for $J_{SS}$ \eqref{eq:triangle_obstruction} (Claim~6.1) &
\textbf{(c)} & \texttt{oneloop\_defect\_phase\_and\_fusion.py} &
Tested on $12$ KP triangles ($36$ edges): $I_{SD}$ fusion $4.51\times 10^{-36}$, triangle closure $1.50\times 10^{-35}$ \\
\addlinespace
Small-amplitude $O(\alpha\ln\alpha)+O(\alpha)$ law \eqref{eq:ISD_small_alpha}, finiteness at absorption lines, and $\pm\pi$ stability/cusp jumps &
\textbf{(c)} & \texttt{oneloop\_defect\_phase\_and\_fusion.py} &
Small-$\alpha$ residual $5.33\times 10^{-19}$ at $\alpha=10^{-6}$; jump errors $<10^{-7}$ \\
\bottomrule
\end{tabular}
\caption{Summary of claims and verification checks in Phase~2g (Route E, step 2).}
\label{tab:summary}
\end{table}

\subsection*{Failures recorded during the investigation}

In accordance with the project ground rules (\texttt{AGENTS.md}), we record all
failed checks and algebraic errors encountered and corrected during this task:
\begin{enumerate}
\item \textbf{Equal weights $(+1,+1)$ instead of dipole weights $(-\eta,+\eta)$ in the 2-root decomposition of $\theta_D'(q;s,\eta)$.}
When first testing the 2-root decomposition of $\theta_D'(q;s,\eta)$ at the
angles $\psi_{1,2} = \gamma \pm \pi/3$, we assigned equal weights $(+1,+1)$ to
$\Theta'(q;\psi_1)$ and $\Theta'(q;\psi_2)$, which produced $O(1)$ residuals
($0.31$). Carrying out the exact partial-fraction decomposition of
$\theta_D'(q;s,\eta)$ showed that the two roots enter with opposite dipole weights
$(c_{D,1},c_{D,2}) = (-\eta,+\eta)$ at $\psi_1 = 2\pi/3-\gamma$ and
$\psi_2 = -2\pi/3-\gamma$, reducing the residual to $1.88\times 10^{-36}$.
\item \textbf{Sign error in the symbolic $4\times 4$ localised-mode determinant.}
In our initial symbolic transcription of the $4\times 4$ linearised CZ sewing
matrix for localised modes, we flipped the sign of the term
$\frac{\sqrt{3}}{6}\eta\sigma(\sigma^2-1)(\kappa_1-\kappa_2)$ in two rows,
which produced an irreducible quintic polynomial in $S = \kappa_1+\kappa_2$ giving
$\gamma_{\mathrm{grow}} = -0.026$ instead of $+0.02863$ at $\eta\sigma = -0.95$.
Restoring the exact sign from row~2 and row~4 of the CZ sewing matrix caused the
determinant to factor completely into $-\frac{S^2(S^2-1)}{24}P_-(S;\eta\sigma)P_+(S;\eta\sigma)$,
matching the numerical $4\times 4$ determinant to $3.89\times 10^{-36}$.
\item \textbf{Spurious $1/\sqrt{3}$ factor in $\sqrt{1-3q_D^2}$ in Check~3 of \texttt{oneloop\_defect\_scattering\_and\_stability.py}.}
When checking $\sin\psi_1 = q_+ - q_D$ for $q_D = -(\cos\gamma)/\sqrt{3}$, we
initially wrote $\sqrt{1-3q_D^2} = (\sin\gamma)/\sqrt{3}$ instead of $\sin\gamma$,
giving a mismatch of $0.207$. Removing the spurious $1/\sqrt{3}$ reduced the
residual to $1.88\times 10^{-36}$.
\item \textbf{Branch cut in $I_{SD}(\alpha,-1;s,\eta)$ and $(\Delta E_D^{\mathrm{ren}},\Delta P_D)$ for $|s|>2$.}
Evaluating $I_{SD}(\alpha,-1;s,\eta)$ by sending $\phi_r \mapsto \pi-\phi_r$
while keeping $\tan(\psi_{s'}/2)^{-1}$ unflipped shifted the integral by $\pi$
for $s<-1$; using the exact crossing relation $\tan((\pi-\phi_r)/2) = \cot(\phi_r/2)$,
which replaces $\tan(\psi_{s'}/2)^{-1}$ by $\tan(\psi_{s'}/2)^\epsilon$ and
multiplies by $\epsilon$, gave agreement to $4.51\times 10^{-36}$ for all
$s \in \mathbb{R}$. Similarly, writing $\ln|\tan(\psi_{s'}/2)|$ instead of the
analytic function $\frac{1}{2}\ln\frac{1-\cos\psi_{s'}}{1+\cos\psi_{s'}}$ in
$\Delta P_D$ and $\Delta E_D^{\mathrm{ren}}$ introduced a $0.2073$ branch-cut
error for $|s|>2$ ($\gamma = \pi+\ii\beta$); using $\frac{1}{2}\ln\frac{1-\cos\psi_{s'}}{1+\cos\psi_{s'}}$
gave agreement to $7.52\times 10^{-37}$.
\item \textbf{Incomplete $\psi$-derivative of $\Theta(q;\psi)$ and complex logarithm in tree-level fusion check.}
In the first run of \texttt{oneloop\_defect\_phase\_and\_fusion.py}, an
over-simplified expression for $\partial_\psi\Theta(q;\psi)$ gave a residual of
$2.1$ in the Faddeev--Korepin $\alpha$-gradient check, and using $\ln(A_{12})$
instead of $\ln|A_{12}|$ in $\delta_0$ raised a \texttt{TypeError} when $A_{12}<0$.
Furthermore, comparing $I_{SD}(3;q_{i+1},\eta)-I_{SD}(3;q_i,\eta)$ directly
against $\rho_{\mathrm{tree},i} I_{SS}(3,S_i)$ without accounting for the
directional transmission factor $\rho_{\mathrm{tree},i} = \epsilon_3\rho_i$ gave
a mismatch for $\epsilon_3=-1$. Fixing $\partial_\psi\Theta(q;\psi)$, using
$\ln|A_{12}|$ and $\ln|z_t|$, and distinguishing $\rho_i = \eta\sgn(q_i-q_{i+1})$
from $\rho_{\mathrm{tree},i} = \epsilon_3\rho_i$ reduced all residuals to
$< 1.67\times 10^{-33}$.
\end{enumerate}

% ==============================================================================
\begin{thebibliography}{9}

\bibitem{CZ2023}
E.~Corrigan and C.~Zambon,
\emph{Integrable defects in the Boussinesq equations},
J.~Phys.~A \textbf{56} (2023) 405202,
\href{https://arxiv.org/abs/2306.03772}{arXiv:2306.03772} [hep-th].
(Full text read in \texttt{literature/2306.03772v2.pdf} / verified via
\texttt{corrigan-zambon/verify\_defect\_solutions.py}.)

\bibitem{FK1978}
L.~D.~Faddeev and V.~E.~Korepin,
\emph{Quantum theory of solitons},
Phys.~Rep.\ \textbf{42} (1978) 1--87,
DOI: \href{https://doi.org/10.1016/0370-1573(78)90058-3}{10.1016/0370-1573(78)90058-3}.
(Full text read in \texttt{literature/Faddeev\_Korepin\_1978\_Quantum\_theory\_of\_solitons.pdf}.)

\bibitem{DHN1974}
R.~F.~Dashen, B.~Hasslacher and A.~Neveu,
\emph{Nonperturbative methods and extended-hadron models in field theory. I. Semiclassical functional methods},
Phys.~Rev.~D \textbf{10} (1974) 4114--4129;
\emph{II. Two-dimensional models and extended hadrons},
Phys.~Rev.~D \textbf{10} (1974) 4130--4138,
DOI: \href{https://doi.org/10.1103/PhysRevD.10.4130}{10.1103/PhysRevD.10.4130}.
(Cited via the full text of \cite{FK1978}, Section~3; original full text not in \texttt{literature/}.)

\bibitem{Hirota1973}
R.~Hirota,
\emph{Exact $N$-soliton solutions of the wave equation of long waves in shallow-water and in nonlinear lattices},
J.~Math.~Phys.\ \textbf{14} (1973) 810--814,
DOI: \href{https://doi.org/10.1063/1.1666400}{10.1063/1.1666400}.
(Full text read in \texttt{literature/}.)

\end{thebibliography}

\end{document}
